% Laget automatisk av elevark/lag_elevark.py fra opplegg/effekt-trapp.html. Endringer her blir overskrevet.
\documentclass[11pt]{article}
\usepackage{elevark}
\begin{document}
\elevarkhode{Fysikk 1 · Mekanisk energi}{Effekten din i trappa}{https://fysikklab.no/opplegg/effekt-trapp.html}{Opplegget på nett}

Når du løper opp en trapp, gjør du et arbeid mot tyngdekraften. Hvor stor effekt klarer du å yte i noen sekunder? Er du sterkere enn en hårføner? Alt dere trenger, er en trapp, en stoppeklokke og en linjal.

\section*{Hensikt}

\begin{itemize}
\item Regne ut endringen i potensiell energi og effekten når du går og løper opp en trapp.
\item Skille mellom energi (hvor mye arbeid) og effekt (hvor fort arbeidet gjøres).
\end{itemize}

\section*{Utstyr}

\begin{itemize}
\item En trapp, en stoppeklokke og en linjal eller et målebånd.
\item Massen trenger ingen å si høyt. Bruk din egen masse bare for deg selv, eller regn med 60~kg.
\end{itemize}

\section*{Teorien}

Når du går opp høyden $h$, øker den potensielle energien med $E_{\mathrm{p}} = mgh$. Effekten er arbeid per tid: $P = E_{\mathrm{p}}/t$. 1 hestekraft er omtrent 746~W.

\section*{Framgangsmåte}

\begin{enumerate}
\item Tell trinnene og mål høyden på ett trinn. Mål gjerne flere trinn og ta gjennomsnittet, eller mål hele høyden med et målebånd.
\item Løperen står med begge føttene på gulvet nedenfor første trinn. Tidtakeren sier «klar, ferdig, gå» og starter klokka på «gå».
\item Løp opp så fort det er trygt, gjerne to trinn om gangen. Stopp klokka når begge føttene er på øverste trinn.
\item Gjør det samme én gang til der den samme personen går i vanlig tempo.
\item Bytt roller, og skriv alle løpene inn i tabellen.
\end{enumerate}

Løp bare i en trapp med rekkverk og uten folk, og ikke løp ned igjen.

\section*{Oppgaver}

\begin{oppgaveliste}
\oppgave \fase{gjett}{Gjett}Hvor stor effekt tror du at du kan yte når du løper opp trappa? Mer eller mindre enn en lyspære (10~W), en hårføner (2000~W) eller en hest (746~W)?
\svar
\oppgave \fase{test}{Mål}Gjør målingene, eller trykk på «Fyll inn eksempel». Hva var den største effekten i gruppa?
\svar
\oppgave \fase{forklar}{Forklar}Sammenlign løpet og gåturen til samme person. Hvorfor er energien den samme, mens effekten er forskjellig?
\svar
\oppgave \fase{test}{Regn}Kroppen har en virkningsgrad på omtrent 25~\% når den gjør arbeid med musklene. Hvor mye kjemisk energi brukte personen i rad 1? Hvor mange ganger må personen gå opp trappa for å bruke energien i en sjokolade på 250 kcal (1 kcal = 4184~J)?
\svar
\oppgave \fase{forklar}{Vurder}Hvilke energiformer og hvilket arbeid har dere sett bort fra? Hvor stor betydning har reaksjonstiden til tidtakeren?
\svar
\oppgave \fase{utfordring}{Utfordring}Hvor lang tid ville en person på 62~kg bruke på å gå opp 500 høydemeter hvis personen kunne holde 600~W hele veien? Hvorfor tar det mye lenger tid i virkeligheten?
\svar
\end{oppgaveliste}

\end{document}
